\documentclass[11pt]{article}

\usepackage[margin=1in]{geometry}
\usepackage{amsmath,amssymb,amsthm,mathtools}
\usepackage{hyperref}
\usepackage{microtype}

\newtheorem{theorem}{Theorem}[section]
\newtheorem{corollary}[theorem]{Corollary}
\newtheorem{proposition}[theorem]{Proposition}
\newtheorem{lemma}[theorem]{Lemma}
\theoremstyle{definition}
\newtheorem{definition}[theorem]{Definition}
\theoremstyle{remark}
\newtheorem{remark}[theorem]{Remark}

\newcommand{\R}{\mathbb R}
\newcommand{\K}{\mathcal K}
\newcommand{\ip}[2]{\langle #1,#2\rangle}
\newcommand{\rank}{\operatorname{rank}}
\newcommand{\spec}{\operatorname{Spec}}
\newcommand{\adj}{\operatorname{adj}}

\title{Boundary-Channel Completion of Rank-Two Displacement Clocks\\
and Finite Zeta Spectral Triples}
\author{}
\date{Draft of July 22, 2026}

\begin{document}
\maketitle

\begin{abstract}
Connes, Consani, and Moscovici construct finite zeta spectral triples from truncated Weil forms
under the assumption that the lowest eigenvalue is simple and its eigenvector is even.  We prove
that neither hypothesis is needed for finite self-adjoint completion.  The result is a universal
boundary-channel trichotomy for a positive semidefinite matrix $T$, a self-adjoint clock $D$, and
the rank-two displacement identity
\[
 DT-TD=|\beta\rangle\langle\eta|-|\eta\rangle\langle\beta|.
\]
If the radical of $T$ meets the even boundary functional, then
$D-|De\rangle\langle\eta|$ preserves the full radical and is self-adjoint in the quotient metric
defined by $T$.  If it meets the odd boundary functional, the dual correction
$D-|Do\rangle\langle\beta|$ has the same properties.  If both functionals vanish on the radical,
a Moore--Penrose source solution gives the neutral correction
$D+|T^\dagger\beta\rangle\langle\eta|$.  The active even and odd cases are mutually exclusive.

For a one-dimensional radical, the neutral case is impossible for the CCM clock.  The even branch
recovers the published real-rooted ground-state transform, while the odd branch produces a new
real-rooted weighted transform
\[
 \sum_j \beta_j\xi_j\prod_{k\ne j}(k-z).
\]
Thus every finite truncated Weil matrix, after subtraction of its spectral floor, canonically
produces a self-adjoint radical quotient without a simplicity or parity assumption.  The theorem
closes the finite completion obstruction; it does not establish the separate cofinal convergence
of normalized determinants to the Riemann $\Xi$ function.  The matrix identities have been
formalized in Lean~4 without additional axioms.
\end{abstract}

\section{Introduction}

The finite spectral construction in \emph{Zeta Spectral Triples} starts from the restriction
$Q_{\lambda}^{N}$ of the Weil quadratic form to the Fourier modes indexed by
$\{-N,\ldots,N\}$~\cite{CCM2025}.  Its matrix is real symmetric, commutes with reflection, and has
the divided-difference form
\[
 q_{ii}=a_i,\qquad q_{ij}=\frac{b_i-b_j}{i-j}\quad(i\ne j),
 \qquad a_{-i}=a_i,\quad b_{-i}=-b_i.                       \tag{1.1}
\]
Writing
\[
 De_j=je_j,\qquad \eta=\sum_{j=-N}^{N}e_j,
 \qquad \beta=\sum_{j=-N}^{N}b_je_j,                       \tag{1.2}
\]
one obtains the exact displacement identity
\[
 DQ_{\lambda}^{N}-Q_{\lambda}^{N}D
 =|\beta\rangle\langle\eta|-|\eta\rangle\langle\beta|.    \tag{1.3}
\]

Let $\epsilon_{\lambda,N}$ be the lowest eigenvalue and set
\[
 T=Q_{\lambda}^{N}-\epsilon_{\lambda,N}I.                  \tag{1.4}
\]
Then $T\geq0$, its radical is the full ground eigenspace, and (1.3) is unchanged.  The construction
in~\cite{CCM2025} assumes that this radical is generated by an even vector $\xi$ normalized by
$\ip{\eta}{\xi}=1$.  The corrected clock
\[
 D'=D-|D\xi\rangle\langle\eta|                              \tag{1.5}
\]
then descends to a self-adjoint operator on the quotient by $\R\xi$.  Its characteristic
determinant is the numerator of the Fourier transform of $\xi$.  The same assumption underlies the
finite real-zero theorem of Connes and van Suijlekom~\cite{CvS2025}.

The simple-even condition is not a consequence of divided-difference displacement: positive
matrices of the same class can have odd or multiple radicals~\cite[Appendices A--B]{CvS2025}.
This makes a source-specific proof of simple-evenness a genuine problem.  Our result shows that it
is not a finite self-adjointness obstruction.  Identity (1.3) already contains a second, parity-dual
boundary channel.  Once all radical vectors are retained, it also contains a neutral source
completion.

The central formula is the radical transport law
\[
 TDv=-\beta\ip{\eta}{v}+\eta\ip{\beta}{v},\qquad v\in\ker T. \tag{1.6}
\]
If an even radical vector has nonzero $\eta$-pairing, (1.6) puts $\beta$ in the range of $T$ and
therefore makes $\beta$ orthogonal to the entire radical.  If an odd radical vector has nonzero
$\beta$-pairing, the same argument puts $\eta$ in the range of $T$.  These alternatives are
mutually exclusive.  The appropriate rank-one correction cancels (1.3) exactly and preserves all
of $\ker T$, not only the chosen vector.

Section~\ref{sec:abstract} proves the abstract trichotomy.  Section~\ref{sec:realzero} derives the
parity-dual real-zero law for a simple radical.  Section~\ref{sec:ccm} applies the theorem to every
finite Weil truncation and separates the closed finite problem from the still-open cofinal
determinant transfer.  Section~\ref{sec:lean} records the formal certificate.

\section{Rank-two displacement systems}
\label{sec:abstract}

Let $V$ be a finite-dimensional real Hilbert space.  All adjoints in this section refer to its
given Euclidean inner product.

\begin{definition}
A \emph{reflection-graded displacement system} is a tuple
$(T,D,\Gamma,\eta,\beta)$ satisfying
\begin{align}
 &T=T^*,\quad D=D^*,\quad T\geq0,\quad \Gamma^2=I,\\
 &\Gamma T=T\Gamma,\quad \Gamma D=-D\Gamma,
 \quad \Gamma\eta=\eta,\quad \Gamma\beta=-\beta,            \tag{2.1}\\
 &DT-TD=|\beta\rangle\langle\eta|-|\eta\rangle\langle\beta|. \tag{2.2}
\end{align}
Put $\K=\ker T$ and $\K_\pm=\ker(\Gamma\mp I)\cap\K$.
\end{definition}

Since $T\geq0$, the formula
\[
 \ip{[x]}{[y]}_T:=\ip{Tx}{y}                               \tag{2.3}
\]
defines a positive-definite inner product on $V/\K$.

\begin{lemma}[Radical transport]
\label{lem:transport}
For every $v\in\K$,
\[
 TDv=-\beta\ip{\eta}{v}+\eta\ip{\beta}{v}.                \tag{2.4}
\]
\end{lemma}

\begin{proof}
Apply (2.2) to $v$ and use $Tv=0$.
\end{proof}

\begin{lemma}[Boundary exclusivity]
\label{lem:exclusive}
The following statements hold.
\begin{enumerate}
\item If $e\in\K_+$ and $\ip{\eta}{e}=1$, then $TDe=-\beta$ and
      $\beta\perp\K$.
\item If $o\in\K_-$ and $\ip{\beta}{o}=1$, then $TDo=\eta$ and
      $\eta\perp\K$.
\end{enumerate}
In particular, the two displayed boundary activities cannot occur simultaneously.
\end{lemma}

\begin{proof}
Parity gives $\ip{\beta}{e}=0$ and $\ip{\eta}{o}=0$, so the two source identities follow from
Lemma~\ref{lem:transport}.  Since $T=T^*$ in finite dimension,
$\operatorname{ran}T=(\ker T)^\perp$.  Thus $\beta=-TDe$ or $\eta=TDo$ is orthogonal to $\K$.
\end{proof}

\begin{theorem}[Boundary-channel completion trichotomy]
\label{thm:trichotomy}
Exactly one of the following three cases holds.
\begin{enumerate}
\item[\rm(E)] $\eta|_{\K}\ne0$.  Choose $e\in\K_+$ with $\ip{\eta}{e}=1$ and set
\[
 A_E=D-|De\rangle\langle\eta|.                             \tag{2.5}
\]
Then $A_E\K\subseteq\K$, $A_Ee=0$, and
\[
 TA_E=A_E^*T.                                               \tag{2.6}
\]
\item[\rm(O)] $\beta|_{\K}\ne0$.  Choose $o\in\K_-$ with $\ip{\beta}{o}=1$ and set
\[
 A_O=D-|Do\rangle\langle\beta|.                            \tag{2.7}
\]
Then $A_O\K\subseteq\K$, $A_Oo=0$, and
\[
 TA_O=A_O^*T.                                               \tag{2.8}
\]
\item[\rm(N)] $\eta|_{\K}=\beta|_{\K}=0$.  Let $y=T^\dagger\beta$ be the unique solution of
$Ty=\beta$ in $\K^\perp$, and set
\[
 A_N=D+|y\rangle\langle\eta|.                              \tag{2.9}
\]
Then $A_N\K\subseteq\K$ and
\[
 TA_N=A_N^*T.                                               \tag{2.10}
\]
\end{enumerate}
Each $A_\bullet$ induces a self-adjoint operator $\overline A_\bullet$ on the Hilbert space
$(V/\K,\ip{\cdot}{\cdot}_T)$; hence its characteristic polynomial has only real zeros.
\end{theorem}

\begin{proof}
The restrictions of $\eta$ and $\beta$ to $\K$ see only $\K_+$ and $\K_-$, respectively.
Lemma~\ref{lem:exclusive} makes (E) and (O) mutually exclusive; if neither occurs, (N) holds.

In case (E), Lemma~\ref{lem:exclusive} and (2.4) give, for $v\in\K$,
\[
 TA_Ev=TDv-TDe\ip{\eta}{v}
       =-\beta\ip{\eta}{v}+\beta\ip{\eta}{v}=0.           \tag{2.11}
\]
Moreover,
\begin{align*}
 TA_E-A_E^*T
 &=(TD-DT)-|TDe\rangle\langle\eta|+|\eta\rangle\langle TDe|\\
 &=\bigl(|\eta\rangle\langle\beta|-|\beta\rangle\langle\eta|\bigr)
   +|\beta\rangle\langle\eta|-|\eta\rangle\langle\beta|=0.
\end{align*}
The normalization gives $A_Ee=0$.  The odd proof is the same calculation with $TDo=\eta$:
\[
 TA_O-A_O^*T
 =(TD-DT)-|\eta\rangle\langle\beta|+|\beta\rangle\langle\eta|=0. \tag{2.12}
\]

In case (N), finite-dimensional self-adjointness gives
$\eta,\beta\in\K^\perp=\operatorname{ran}T$, so $y$ exists.  Equation (2.4) shows
$D\K\subseteq\K$, while $\eta|_\K=0$ makes the correction vanish on $\K$.  Finally,
\[
 TA_N-A_N^*T
 =(TD-DT)+|\beta\rangle\langle\eta|-|\eta\rangle\langle\beta|=0. \tag{2.13}
\]

In every case, radical preservation makes the quotient operator well defined.  For quotient
classes,
\[
 \ip{\overline A[x]}{[y]}_T
 =\ip{TAx}{y}=\ip{Tx}{Ay}=\ip{[x]}{\overline A[y]}_T,
\]
so it is self-adjoint.
\end{proof}

\begin{remark}[Affine neutral completions]
The neutral branch is not tied to the displayed pseudoinverse formula.  If $Tx=\eta$ and
$Ty=\beta$, both $D+|y\rangle\langle\eta|$ and
$D-|x\rangle\langle\beta|$ complete the metric adjointness defect.  Convex affine combinations
do as well.  The Moore--Penrose choice is canonical because it lies in $\K^\perp$.
\end{remark}

\section{The parity-dual real-zero law}
\label{sec:realzero}

Assume now that $\K=\R\xi$ and that $D$ has simple real spectrum.  In the CCM clock,
$\ker D=\R e_0$, where $\Gamma e_0=e_0$ and $\ip{\eta}{e_0}=1$.

\begin{lemma}
\label{lem:noneutral}
The neutral branch is impossible.  Hence a simple radical is active even or active odd.
\end{lemma}

\begin{proof}
Reflection invariance makes $\xi$ even or odd.  If both boundary pairings vanished,
Lemma~\ref{lem:transport} would give $D\xi\in\K$.  Since $D$ reverses parity while $\K$ has one
parity, $D\xi=0$.  Thus $\xi$ is proportional to $e_0$, contradicting
$\ip{\eta}{e_0}=1$.
\end{proof}

\begin{theorem}[Parity-dual real-zero determinants]
\label{thm:realzero}
Let $K=\R\xi$.
\begin{enumerate}
\item If $\xi$ is even and $\ip{\eta}{\xi}=1$, then
\[
 P_E(z)=\det(D-zI)\ip{\eta}{(D-zI)^{-1}\xi}                \tag{3.1}
\]
is the characteristic polynomial of $\overline A_E$ and has only real zeros.
\item If $\xi$ is odd and $\ip{\beta}{\xi}=1$, then
\[
 P_O(z)=\det(D-zI)\ip{\beta}{(D-zI)^{-1}\xi}              \tag{3.2}
\]
is the characteristic polynomial of $\overline A_O$ and has only real zeros.
\end{enumerate}
Both identities extend polynomially through $\spec D$.
\end{theorem}

\begin{proof}
Let $c=\eta$ in the even case and $c=\beta$ in the odd case.  The correction has the common form
$A=D-|D\xi\rangle\langle c|$, with $\ip{c}{\xi}=1$.  Away from $\spec D$, the matrix determinant
lemma gives
\begin{align*}
 \det(A-zI)
 &=\det(D-zI)\left(1-\ip{c}{(D-zI)^{-1}D\xi}\right)\\
 &=-z\det(D-zI)\ip{c}{(D-zI)^{-1}\xi}.                    \tag{3.3}
\end{align*}
Since $A\xi=0$ and $A$ preserves $\R\xi$, the left side is
$-z\det(\overline A-zI)$.  Canceling $-z$ proves (3.1)--(3.2).  The quotient operator is
self-adjoint by Theorem~\ref{thm:trichotomy}.
\end{proof}

If $De_j=je_j$, formula (3.2) becomes
\[
 P_O(z)=\sum_{j=-N}^{N}\beta_j\xi_j
          \prod_{\substack{k=-N\\k\ne j}}^{N}(k-z).       \tag{3.4}
\]
The even formula has $\eta_j=1$ and is the polynomial appearing in the CCM Fourier transform.
The odd formula is its displacement-weighted dual.

\section{Application to finite zeta spectral triples}
\label{sec:ccm}

\begin{corollary}[Unconditional finite completion]
For every aperture $\lambda>1$ and every Fourier truncation $N$, let
$\epsilon_{\lambda,N}$ be the lowest eigenvalue of the actual truncated Weil matrix and put
$T=Q_\lambda^N-\epsilon_{\lambda,N}I$.  Then one of the completions in
Theorem~\ref{thm:trichotomy} induces a self-adjoint clock on the quotient by the full ground
eigenspace.  No assumption on ground multiplicity or parity is required.
\end{corollary}

\begin{proof}
The matrix $T$ is positive semidefinite and retains the displacement identity (1.3), because a
scalar matrix commutes with $D$.  Reflection properties follow from (1.1).  Apply
Theorem~\ref{thm:trichotomy}.
\end{proof}

This corollary changes the logical boundary of the CCM program.  The finite self-adjoint quotient
does not require a proof that the actual ground is simple even.  What simple-evenness supplies is
a specific identification: the boundary determinant is then the unweighted Fourier transform of
the even ground vector, which is the branch compared in~\cite{CCM2025} with the prolate
$h_0-h_4$ seed.  An odd ground produces the equally real-rooted polynomial (3.4), but no theorem
currently identifies its cofinal limit with $\Xi$.

Nor does finite completion prove the second missing CCM step.  To deduce RH from finite real
spectra one still needs normalization and locally uniform convergence strong enough to transport
the divisor to $\Xi$.  The present theorem supplies no such estimate.  Its contribution is
orthogonal and exact: it removes the parity/multiplicity obstruction from every finite stage and
exposes cofinal branch selection as the remaining source problem.

\section{Formal verification}
\label{sec:lean}

The algebraic core is formalized in the Lean~4 module
\texttt{GaloisForLFunctions.RHTower.CCMBoundaryChannelCompletion}.  The certificate proves:
\begin{itemize}
\item the radical transport identity (2.4);
\item even/odd boundary exclusivity;
\item all three metric-adjointness identities $TA=A^\mathsf{T}T$;
\item annihilation of the normalized active vector;
\item preservation of every vector in the full radical;
\item bundled active-even and active-odd connector theorems from displacement alone.
\end{itemize}
The declarations use only Lean's standard logical axioms
\texttt{propext}, \texttt{Classical.choice}, and \texttt{Quot.sound}; no mathematical axiom,
\texttt{sorry}, or RH-facing hypothesis is introduced.  The quotient spectral theorem and the
determinant lemma are standard finite-dimensional consumers of the certified identities; their
specialized polynomial packaging is left for a subsequent formal layer.

\section{Contribution and open boundary}

The universal law is the source/radical classification
\[
\begin{array}{ccl}
\eta|_\K\ne0 &\Longrightarrow& \eta\text{-boundary completion and }\beta\perp\K,\\
\beta|_\K\ne0&\Longrightarrow& \beta\text{-boundary completion and }\eta\perp\K,\\
\eta|_\K=\beta|_\K=0&\Longrightarrow&\text{neutral source-solution completion}.
\end{array}                                                     \tag{6.1}
\]
It applies beyond the zeta setting to every reflection-graded rank-two displacement pencil.  Its
finite CCM consequence is unconditional and solves the need for a simple-even hypothesis in the
construction of a self-adjoint radical quotient.

Two problems remain separate.  First, the actual source may still select a simple even ground;
our theorem neither proves nor needs that assertion at a finite stage.  Second, a proof of the
Riemann hypothesis through this lane still requires cofinal source/divisor transport from a
specified branch to $\Xi$.  Finite real-rootedness alone cannot provide it.  Any future transfer
theorem must therefore do more than re-establish finite self-adjointness: it must identify the
active channel and control its normalized determinant in the cofinal aperture limit.

\begin{thebibliography}{9}

\bibitem{CCM2025}
A.~Connes, C.~Consani, and H.~Moscovici,
\emph{Zeta Spectral Triples}, arXiv:2511.22755 (2025),
\url{https://arxiv.org/abs/2511.22755}.

\bibitem{CvS2025}
A.~Connes and W.~D. van Suijlekom,
\emph{Quadratic Forms, Real Zeros and Echoes of the Spectral Action},
arXiv:2511.23257 (2025),
\url{https://arxiv.org/abs/2511.23257}.

\end{thebibliography}

\end{document}
